\documentclass[11pt,a4paper]{article}
\usepackage[utf8]{inputenc}
\usepackage[T1]{fontenc}
\usepackage[brazilian]{babel}
\usepackage{geometry}
\geometry{top=1.5cm, bottom=1.5cm, left=1.6cm, right=1.6cm}
\usepackage{xcolor}
\usepackage{titlesec}
\usepackage{enumitem}
\usepackage{booktabs}
\usepackage{tabularx}
\usepackage{tcolorbox}
\usepackage{hyperref}
\usepackage{microtype}
\usepackage{newtxtext,newtxmath}

% Paleta Institucional IFES
\definecolor{ifesgreen}{RGB}{24, 138, 63}
\definecolor{ifesdark}{RGB}{14, 88, 39}
\definecolor{ifesred}{RGB}{200, 30, 30}
\definecolor{darkgray}{RGB}{40, 44, 52}
\definecolor{lightgray}{RGB}{248, 249, 250}
\definecolor{speechbg}{RGB}{240, 248, 242}
\definecolor{bordergreen}{RGB}{180, 220, 195}

\hypersetup{
    colorlinks=true,
    linkcolor=ifesgreen,
    urlcolor=ifesgreen,
    citecolor=ifesgreen
}

\titleformat{\section}{\color{ifesgreen}\large\bfseries}{\thesection}{1em}{}[\color{ifesgreen}\titlerule]

\tcbset{
    boxrule=0.8pt,
    arc=4pt,
    left=6pt, right=6pt, top=4pt, bottom=4pt
}

\begin{document}

% Cabeçalho
\begin{center}
    {\small\textbf{INSTITUTO FEDERAL DO ESPÍRITO SANTO -- IFES | CAMPUS SERRA}}\\[2pt]
    {\small\textbf{PROGRAMA DE PÓS-GRADUAÇÃO EM ENGENHARIA DE CONTROLE E AUTOMAÇÃO (PROPECAUT)}}\\[4pt]
    \rule{\textwidth}{1.2pt}\\[5pt]
    {\large\textbf{\color{ifesgreen}ROTEIRO DE FALA E GUIA DE APRESENTAÇÃO (8 MINUTOS)}}\\[2pt]
    {\small\textbf{Dissertação de Mestrado: Controle e Gerenciamento Térmico-Energético de Data Center Orbital}}\\[4pt]
    \rule{\textwidth}{0.6pt}
\end{center}

\vspace{-6pt}
\begin{center}
    \small
    \begin{tabularx}{\textwidth}{lXlX}
        \textbf{Mestrando:} & Dafnhê Thor C. M. Barros & \textbf{Tempo Alvo:} & \textbf{8 minutos (480 s)} \\
        \textbf{Orientador:} & Prof. Dr. Flávio B. Braz da Silva & \textbf{Slides:} & 7 telas no PPTX Widescreen
    \end{tabularx}
\end{center}

\vspace{2pt}

\begin{tcolorbox}[colback=lightgray, colframe=ifesgreen, title=\color{ifesgreen}\textbf{DIRETRIZ DE POSTURA E TEMPO PARA APRESENTAÇÃO DE 8 MINUTOS}]
\small
\textbf{Ritmo de fala:} Mantenha fala firme, pausada e sem pressa. Em 8 minutos, você não deve tentar explicar deduções matemáticas de 30 linhas; o seu objetivo é transmitir com clareza cristalina: \textbf{(1) A dor física real}, \textbf{(2) A lacuna científica}, \textbf{(3) O coração da sua solução em Controle (PINN + TEMS)} e \textbf{(4) Os resultados e viabilidade do modelo}.
\end{tcolorbox}

\vspace{6pt}

% SLIDE 1
\begin{tcolorbox}[colback=speechbg, colframe=bordergreen, title=\color{ifesdark}\textbf{SLIDE 1: CAPA OFICIAL \hfill TEMPO: 0:00 -- 0:30 (Duração: 30 segundos)}]
\small
\textbf{Objetivo do slide:} Apresentar-se formalmente, enunciar o título exato e situar a pesquisa no programa.\\
\textbf{Texto sugerido para falar:}
\begin{quote}
``Bom dia a todos os membros da banca examinadora e aos presentes. Meu nome é Dafnhê Thor Castro de Martins Barros, sou orientado pelo Professor Doutor Flávio Barcelos, e hoje apresento o projeto de dissertação intitulado: \textbf{Controle e Gerenciamento Térmico-Energético Computacional de um Modelo de Data Center Orbital}, desenvolvido no Programa de Pós-Graduação em Engenharia de Controle e Automação do IFES Campus Serra.''
\end{quote}
\end{tcolorbox}

\vspace{6pt}

% SLIDE 2
\begin{tcolorbox}[colback=speechbg, colframe=bordergreen, title=\color{ifesdark}\textbf{SLIDE 2: CONTEXTO, OPORTUNIDADE ORBITAL E A LACUNA \hfill TEMPO: 0:30 -- 1:45 (Duração: 75 segundos)}]
\small
\textbf{Objetivo do slide:} Criar a ponte entre o problema terrestre de IA e a oportunidade no espaço, apontando a lacuna de controle.\\
\textbf{Texto sugerido para falar:}
\begin{quote}
``Para contextualizar: a explosão da inteligência artificial generativa criou uma crise sem precedentes na infraestrutura terrestre. Data centers já consomem parcelas gigantescas das redes elétricas das cidades e utilizam bilhões de litros de água tratada para arrefecimento evaporativo. 

A indústria aeroespacial propõe levar essas cargas para o espaço. Em órbita heliossíncrona, a radiação solar AM0 é ininterrupta, fornecendo até oito vezes mais energia que usinas em solo, com consumo de água zero. 

Porém, surge um gargalo físico severo: no vácuo espacial não existe ar, e portanto inexiste convecção ou condução com o ambiente. Todo o calor gerado pelos processadores de IA deve ser rejeitado para o espaço profundo estritamente por radiação térmica. A literatura atual trata essa potência de forma estática; a lacuna que preenchemos é a criação de um framework dinâmico que realize o \textbf{controle térmico-energético ativo em tempo real} para evitar o colapso dos chips.''
\end{quote}
\end{tcolorbox}

\vspace{6pt}

% SLIDE 3
\begin{tcolorbox}[colback=speechbg, colframe=bordergreen, title=\color{ifesdark}\textbf{SLIDE 3: OBJETIVOS E METODOLOGIA COMPUTACIONAL \hfill TEMPO: 1:45 -- 3:00 (Duração: 75 segundos)}]
\small
\textbf{Objetivo do slide:} Definir o objetivo geral, a escala adotada (50 kW) e defender a abordagem estritamente computacional.\\
\textbf{Texto sugerido para falar:}
\begin{quote}
``O objetivo geral desta dissertação é desenvolver, parametrizar e validar computacionalmente um ambiente integrado para o gerenciamento térmico-energético de um modelo de data center orbital de 50 kW úteis de processamento de IA.

Do ponto de vista metodológico, adotamos uma abordagem estritamente \textit{in silico}, baseada em modelagem matemática analítica rigorosa e simulação computacional. Justificamos essa escolha pela impossibilidade física e orçamentária de ensaiar experimentalmente uma carga de 50 kW em câmaras de termovácuo (TVAC) laboratoriais. Todos os nossos parâmetros são fundamentados em normas da NASA, ECSS e datasheets de componentes espaciais qualificados.

Todo o ecossistema foi construído em Python científico, utilizando bibliotecas de ponta como DeepXDE e PyTorch para redes neurais informadas pela física, e mantido sob licença de Ciência Aberta com repositório público no GitHub para plena reprodutibilidade pela banca.''
\end{quote}
\end{tcolorbox}



% SLIDE 4
\begin{tcolorbox}[colback=speechbg, colframe=bordergreen, title=\color{ifesdark}\textbf{SLIDE 4: FÍSICA DA DISSIPAÇÃO NO VÁCUO E O LIMITE DOS CHIPS \hfill TEMPO: 3:00 -- 4:15 (Duração: 75 segundos)}]
\small
\textbf{Objetivo do slide:} Explicar a barreira de Stefan-Boltzmann e demonstrar por que o controle precisa impedir a fuga térmica.\\
\textbf{Texto sugerido para falar:}
\begin{quote}
``Entrando na física do problema: a taxa de rejeição radiativa é governada pela Lei de Stefan-Boltzmann, proporcional à quarta potência da temperatura. A uma temperatura típica de radiador de 70 °C, conseguimos dissipar no máximo cerca de 708 Watts por metro quadrado. Isso significa que cada quilowatt elétrico exige mais de 1,4 metros quadrados de painéis radiadores. Por essa razão, adotamos uma arquitetura modular em cluster com nós de 50 kW, contendo a área radiativa em 88,6 metros quadrados por satélite e evitando estruturas excessivamente flexíveis que desestabilizariam o controle de atitude.

Do lado dos semicondutores de IA, a potência consumida possui uma parcela dinâmica que varia com o clock e o quadrado da tensão, e uma corrente de fuga que cresce exponencialmente com a temperatura. Se o calor se acumula, o chip entra em fuga térmica (\textit{thermal runaway}). Portanto, estabelecemos como restrição rígida de projeto manter a temperatura de junção sempre abaixo de 80 °C.''
\end{quote}
\end{tcolorbox}

\vspace{6pt}

% SLIDE 5
\begin{tcolorbox}[colback=speechbg, colframe=bordergreen, title=\color{ifesdark}\textbf{SLIDE 5: A SOLUÇÃO: PREDITOR PINN E CONTROLE TEMS \hfill TEMPO: 4:15 -- 6:00 (Duração: 105 segundos)}]
\small
\textbf{Objetivo do slide:} Apresentar o coração da dissertação em Engenharia de Controle: a PINN rápida e o controlador TEMS.\\
\textbf{Texto sugerido para falar:}
\begin{quote}
``Aqui reside a principal contribuição da nossa pesquisa em Engenharia de Controle. 

Para controlar a temperatura em malha fechada a bordo de um satélite, precisaríamos prever o comportamento térmico com antecedência. Métodos numéricos convencionais, como Elementos Finitos (FEM), levam minutos ou horas para resolver a equação de calor com contorno radiativo não linear, o que é inviável para processadores de bordo espaciais operando entre 100 e 200 MHz.

Para contornar esse gargalo, empregamos uma \textbf{Rede Neural Informada pela Física (PINN)}. A PINN incorpora a equação diferencial de condução e o contorno de Stefan-Boltzmann diretamente na função de perda durante o treino. Uma vez treinada, a sua inferência é uma avaliação matricial direta executada em menos de 5 milissegundos.

A PINN atua, portanto, como o \textbf{modelo substituto preditivo} para o nosso controlador, o \textbf{TEMS} (\textit{Thermal-Energy Management System}). O TEMS lê a predição da rede e atua em duas frentes: modula preventivamente a tensão e o clock dos processadores (via DVFS) e redistribui lotes de tarefas de IA entre nós do cluster, garantindo matematicamente que a barreira de 80 °C jamais seja violada.''
\end{quote}
\end{tcolorbox}

\vspace{6pt}

% SLIDE 6
\begin{tcolorbox}[colback=speechbg, colframe=bordergreen, title=\color{ifesdark}\textbf{SLIDE 6: RESULTADOS ESPERADOS E CONTRIBUIÇÕES \hfill TEMPO: 6:00 -- 7:15 (Duração: 75 segundos)}]
\small
\textbf{Objetivo do slide:} Demonstrar o ganho comparativo do controle e consolidar o fechamento formal.\\
\textbf{Texto sugerido para falar:}
\begin{quote}
``Na validação comparativa do modelo, analisamos três cenários operacionais:
No primeiro cenário, sem controle ativo, uma rajada de processamento leva os chips a ultrapassarem 100 °C em menos de oito minutos, destruindo a pastilha de silício.
No segundo cenário, com controle reativo clássico por corte abrupto (\textit{thermal throttling}), o sistema sofre de oscilações severas e perda drástica de vazão computacional.
No terceiro cenário, aplicando o nosso framework preditivo TEMS acoplado à PINN, a temperatura estabiliza-se suavemente abaixo de 80 °C, mantendo a produtividade computacional máxima com um PUE orbital de aproximadamente 1,05.

Como desdobramento futuro, propomos a evolução desse modelo para um Gêmeo Digital Operacional com calibração contínua por telemetria em tempo real. O trabalho cumpre integralmente os requisitos regimentais do ProPECAut e gerará artigos científicos em periódicos qualificados da área.''
\end{quote}
\end{tcolorbox}

\vspace{6pt}

% SLIDE 7
\begin{tcolorbox}[colback=speechbg, colframe=bordergreen, title=\color{ifesdark}\textbf{SLIDE 7: ENCERRAMENTO E AGRADECIMENTOS \hfill TEMPO: 7:15 -- 8:00 (Duração: 45 segundos)}]
\small
\textbf{Objetivo do slide:} Finalizar com segurança e abrir formalmente para a arguição.\\
\textbf{Texto sugerido para falar:}
\begin{quote}
``Concluo agradecendo ao Instituto Federal do Espírito Santo, ao ProPECAut e ao meu orientador, Prof. Flávio, pelo apoio contínuo. Convido os membros da banca a acessarem os modelos em nosso repositório no GitHub e coloco-me à inteira disposição para as perguntas, comentários e considerações da banca examinadora. Muito obrigado!''
\end{quote}
\end{tcolorbox}

\end{document}
